\section{Methods}
\label{sec:methods}

\subsection{Study Objective}

The objective of this study is to evaluate whether a protocol-structured analysis (``LLM Protocols'' condition) produces explanations of major failures that are closer to independent assessor reports than unstructured baseline large language model (LLM) analyses. 

Closeness is operationalized through blinded expert pairwise comparison using a cumulative structural rubric (R1--R4). The primary hypothesis is conditional:

\begin{quote}
Protocol-structured analyses will more frequently be judged closer to independent assessor reports in governance-heavy, underdetermined cases (R3--R4 dominant), with smaller or no gains in mechanically crisp cases (R1--R2 dominant).
\end{quote}

The study does not claim universal superiority. Instead, it tests whether structured reasoning confers measurable advantage under specific structural conditions, while remaining neutral in others.

\subsection{Design Overview}
This study evaluates the FRFP reasoning protocol as an alternative
agent configuration within a controlled AI-agent evaluation framework.
The experimental design holds the task dataset, model configuration,
and evaluation procedure constant while varying the analytical workflow
applied by the agent.
The study uses a three-condition within-case design across 25 historical major failures.  For each case, three analyses are generated; each corresponding to a distinct AI agent configuration.
\begin{enumerate}
    \item \textbf{Baseline (B):} Unstructured engineering-style analysis generated by an LLM without protocol framing.
    \item \textbf{Protocol (P):} Analysis generated under explicit protocol constraints derived from the Foundational Reasoning and Feedback Protocol (FRFP Vol.~1).
    \item \textbf{Independent (I):} Report-grounded summary constrained to publicly available independent assessor findings.
\end{enumerate}

All three outputs are normalized into structurally identical one-page documents prior to evaluation. Expert raters compare the three blinded documents per case and identify the closest pair, along with a cumulative similarity level (R1--R4).

A key methodological feature of the design is that the primary comparison
between Baseline and Protocol holds information constant while varying only
the analytical workflow. This allows differences in output quality to be
attributed to protocol structure rather than to differential access to
external evidence.

\subsection{Case Selection and Stratification}

To mitigate cherry-picking, cases were selected using a pre-specified inclusion rule:

\begin{itemize}
    \item A major failure with substantial real-world impact (loss of life, financial loss, systemic disruption, or regulatory consequence).
    \item Availability of a publicly accessible independent assessor report (e.g., accident board, regulatory commission, court findings).
    \item Sufficient public documentation of system setup and decision process.
\end{itemize}

Cases were stratified across five solution-space categories:

\begin{enumerate}
    \item Mechanically crisp single-chain failures
    \item Socio-technical failures with governance components
    \item Financial/institutional failures involving metrics and model drift
    \item Multi-actor regulatory/certification breakdowns
    \item Crisis management and disaster response failures
\end{enumerate}

The 25 cases include:

\begin{itemize}
    \item Ariane 5 Flight 501
    \item Space Shuttle Columbia (CAIB)
    \item Space Shuttle Challenger (Rogers Commission)
    \item Boeing 737 MAX (JT610, ET302 reports)
    \item Deepwater Horizon (National Commission)
    \item Fukushima Daiichi
    \item Three Mile Island
    \item Chernobyl
    \item London Whale (U.S. Senate PSI Report)
    \item Barings Bank Collapse
    \item Société Générale / Kerviel
    \item Knight Capital
    \item Therac-25
    \item Mars Climate Orbiter
    \item Mars Polar Lander
    \item Patriot Missile (Dhahran)
    \item Toyota Unintended Acceleration
    \item Exxon Valdez
    \item Hawaii False Missile Alert
    \item Equifax Breach
    \item Heartbleed
    \item Wells Fargo Fake Accounts
    \item Uber Tempe Autonomous Vehicle Crash
    \item Vermont Gas ANGP
\end{itemize}

\subsection{Generation Phase (Isolation Protocol)}

Each case was generated under three strictly separated prompt envelopes:

\paragraph{Baseline Condition (B)}
Prompt: Engineering-style failure analysis without reference to protocol frameworks, gates, or structural formalism.

\paragraph{Protocol Condition (P)}
Prompt: Explicit application of protocol primitives including:
\begin{itemize}
    \item Task definition
    \item Control surfaces
    \item Validation gates
    \item Escalation mechanisms
    \item Observability constraints
\end{itemize}

\paragraph{Independent Condition (I)}
Prompt: Summary constrained to independent assessor findings; no reinterpretation beyond report conclusions.

Generation was performed in separate conversational contexts to avoid cross-contamination of structure or vocabulary.

\subsection{Normalization Phase}

All generated outputs were normalized into a fixed one-page shell:

\begin{enumerate}
    \item System / Task Description
    \item Structural Model of the Process
    \item Observed Outcome
    \item Failure Mechanism
    \item Governance / Control Breakdown
    \item Preventative Interventions
\end{enumerate}

Normalization rules:
\begin{itemize}
    \item No new reasoning introduced.
    \item No vocabulary injection from protocol condition into baseline.
    \item No expansion beyond original content.
    \item Word budget standardized (500--650 words).
\end{itemize}

This ensures formatting uniformity while preserving internal reasoning differences.

\subsection{Blinding and Randomization}

For each case, the three normalized documents were labeled Type X, Type Y, and Type Z. The mapping between document type (B, P, I) and X/Y/Z was randomized per rater.

Each rater received a personalized packet:
\begin{itemize}
    \item 25 cases
    \item Randomized case order
    \item Randomized X/Y/Z mapping per case
\end{itemize}

An administrator key linking packet ID to document type was stored separately and not shared with raters.

\subsection{Participants}

Ten expert raters were recruited based on experience in:
\begin{itemize}
    \item Engineering systems
    \item Risk management
    \item Regulatory oversight
    \item Safety-critical domains
\end{itemize}

Each rater evaluated all 25 cases independently.

\subsection{Evaluation Instrument}

For each case, raters selected:

\begin{enumerate}
    \item The closest pair among (X\&Y, Y\&Z, Z\&X)
    \item A cumulative similarity level:
\end{enumerate}

\begin{itemize}
    \item R1: Control surface match
    \item R2: R1 + primary mechanism match
    \item R3: R2 + decision points/timeline match
    \item R4: R3 + institutional/governance mechanism match
\end{itemize}

Raters were instructed to select the highest level that remained valid.

Optional justification bullets (J1--J5) captured which dimensions influenced their decision.

\subsection{Primary Outcome Measure}

For each case, define:

\[
P_{\text{Protocol}} = \text{Proportion of raters selecting (Protocol, Independent) as closest pair}
\]

\[
P_{\text{Baseline}} = \text{Proportion selecting (Baseline, Independent)}
\]

Primary effect size:
\[
\Delta = P_{\text{Protocol}} - P_{\text{Baseline}}
\]

\subsection{Statistical Analysis}

\begin{itemize}
    \item Per-case comparison using binomial proportions.
    \item Stratified analysis by case category.
    \item Bootstrap confidence intervals across cases.
    \item Secondary analysis restricted to R3--R4 selections.
\end{itemize}

Mechanically crisp cases are analyzed separately as negative controls to evaluate whether protocol structuring introduces degradation, unnecessary complexity, or inflation of unsupported claims in single-chain failures.

\subsection{Robustness Checks}

\begin{itemize}
    \item Verify no formatting bias via normalization constraints.
    \item Conduct subset analysis excluding vocabulary unique to protocol condition.
    \item Report full per-case scorecard in appendix.
\end{itemize}

\subsection{Reproducibility}

All prompts, normalization rules, randomization seeds, and packet mappings are archived. Independent replication is possible using the published protocol and case list.

\subsection{Batch Generation Protocol}
\label{sec:batch-generation}

\paragraph{Objective.}
The generation phase produces three analytically independent documents per case: (i) a Baseline LLM analysis (B), (ii) a Protocol-structured LLM analysis (P), and (iii) an Independent assessor-grounded summary (I). The protocol enforces cognitive isolation across conditions, freezes raw outputs prior to normalization, and maintains an auditable trail sufficient for replication and review.

\subsubsection{Case Freeze}

Prior to generation, the full list of 25 cases is finalized and assigned unique case IDs. For each case, a neutral case description (target length: $\leq 150$ words) is prepared, focusing on system setup and task intent rather than failure diagnosis. The case list and descriptions are frozen (no additions, removals, or wording edits) and archived with a timestamp and cryptographic hash (e.g., SHA-256).

\subsubsection{Model and Environment Fixing}

To minimize generation drift, the model version and decoding parameters are fixed and logged (e.g., temperature, top-$p$, frequency and presence penalties). The system prompt is held constant across conditions except for condition-specific instructions described below. Unless explicitly required for the Independent condition, browsing and external retrieval are disabled. All parameters are archived.

\subsubsection{Three-Pipeline Isolation}

Each condition is generated under strict isolation:
\begin{itemize}
    \item Each condition is generated in a separate session with conversation state reset.
    \item No outputs from any condition are visible during generation of another condition.
    \item Generation order (B/P/I) is randomized per case to reduce ordering effects.
\end{itemize}

\subsubsection{Condition B: Baseline Analysis}

\paragraph{Prompt envelope.}
The Baseline prompt requests a structured post hoc analysis in plain engineering language without invoking any formal reasoning protocol. It instructs the model to remain neutral, avoid protocol terminology, and avoid direct reliance on independent investigation reports.

\paragraph{Constraints.}
\begin{itemize}
    \item No protocol primitives (e.g., explicit/tacit separation, gates) are introduced.
    \item No iterative refinement; single-pass generation only.
    \item Regeneration is allowed only for technical failure or clear constraint violation.
\end{itemize}

\subsubsection{Condition P: Protocol-Structured Analysis}

\paragraph{Prompt envelope.}
The Protocol condition applies FRFP-style structural primitives to produce an auditable process model,
including task intent; representation and measurement layers; control surfaces; escalation/endorsement nodes;
observability constraints; and governance failure points. The model must not cite or quote independent assessor
reports. 

The protocol prompt implements the executable FRFP-compatible
instruction surface defined in the FRFP Technical Specification.

\paragraph{Executable surface (normative).}
Condition P is the only condition that uses the executable FRFP-compatible instruction surface:
(i) intent locking plus ambiguity identification,
(ii) artifact discipline with explicit labels for \texttt{FACT}/\texttt{ASSUMPTION}/\texttt{MECHANISM} (including uncertainty),
(iii) separation of generation and endorsement (non-binding language; no self-endorsement; explicit handoff plan to human evaluation),
and (iv) canonical ordering (structural reconstruction $\rightarrow$ diagnostics/branch exploration $\rightarrow$
handoff/verification paths). Diagnostics and handoff must be explicit in the response (typically as the closing portion of
the output).

\emph{Interpretation note:} the handoff component is produced only as an explicit artifact (a proposed verification/review plan);
no post-generation human endorsement or modification is performed within Condition P for the purposes of the survey comparison.

\paragraph{Constraints.}
\begin{itemize}
  \item Protocol decomposition is mandatory (decision nodes and control points must be explicit).
  \item No browsing or external retrieval.
  \item Single-pass generation only (no iterative refinement).
\end{itemize}

\subsubsection{Condition I: Independent Assessor Summary}

\paragraph{Prompt envelope.}
The Independent condition requests a concise summary constrained to publicly documented independent assessor findings (e.g., investigation boards, commissions, regulators). The prompt prohibits interpretive extrapolation beyond documented conclusions and prohibits importing protocol terminology.

\paragraph{Constraints.}
\begin{itemize}
    \item Grounding is required; if browsing is enabled, citations/notes are recorded in the research log.
    \item No protocol framing or theoretical language is introduced.
    \item Single-pass generation only; regeneration allowed only for technical failure or clear constraint violation.
\end{itemize}

\subsubsection{Raw Output Freezing and Archival}

Immediately after generation, each output is saved as a raw artifact (e.g., \texttt{CaseID\_B\_raw.txt}, \texttt{CaseID\_P\_raw.txt}, \texttt{CaseID\_I\_raw.txt}), hashed (e.g., SHA-256), and timestamped. Raw outputs are not edited. These frozen raw files constitute the primary dataset for subsequent normalization and evaluation.

\subsubsection{Normalization Phase (Post-Freezing)}

Only after all three raw outputs for a case are frozen, a separate normalization pass converts each output into a structurally identical one-page format using a fixed six-section shell:
\begin{enumerate}
    \item System / Task Description
    \item Structural Model of the Process
    \item Observed Outcome
    \item Failure Mechanism
    \item Governance / Control Breakdown
    \item Preventative Interventions
\end{enumerate}

Normalization constraints:
\begin{itemize}
    \item No new reasoning is added and no semantic enrichment is performed.
    \item No vocabulary injection from Protocol into Baseline is permitted.
    \item Presentation is standardized (word budget target: 500--650 words; tone and formatting harmonized).
\end{itemize}

Normalization is performed blind to condition labels. Normalized artifacts are saved as Type X/Y/Z documents for distribution to raters; the mapping from (B,P,I) to (X,Y,Z) is stored separately in an administrative key.

\subsubsection{Contamination and Leakage Checks}

Before packet distribution, an integrity check is performed to detect cross-condition contamination (e.g., anomalous shared phrasing, protocol terminology appearing in Baseline or Independent, or structural leakage beyond the fixed shell). If contamination is detected, the case is regenerated from scratch under the isolation protocol and the discarded versions are logged.

\subsubsection{Rater Packet Randomization and Blinding}

For each rater, case order is randomized and Type X/Y/Z ordering is randomized per case using a fixed seed recorded in the administrative archive. Raters receive only Type labels and are not informed of the existence or nature of the underlying conditions. The administrative key linking packet IDs to condition labels is not accessible to raters during evaluation.

\subsubsection{Regeneration Policy}

Regeneration is permitted only for: (i) technical failure, (ii) incompleteness, or (iii) explicit prompt constraint violation. Regeneration is not permitted due to perceived quality imbalance or stylistic preference. All regeneration events, including rationale and hashes of discarded outputs, are recorded.

\subsubsection{Reproducibility Commitments}

The frozen case list, prompts, model parameters, raw outputs, normalized outputs, randomization seeds, packet mappings, and hash logs are archived. This enables independent replication of the generation, normalization, and evaluation pipeline under the same experimental specification

\subsubsection{Case List Finalization and Freeze}
\label{sec:case-freeze}

The full list of 25 cases is finalized exactly as specified in the FRFPVsBaseSurvey document \cite{frfp_base_survey}. No additions, removals, substitutions, or reclassification of cases are permitted after this point.

Each case is assigned a unique numeric Case ID (C01--C25). The mapping from case name to Case ID is stored in an administrative registry and is not modified after freezing.

\paragraph{Neutral Case Descriptions.}
For each case, a neutral case description (target length: $\leq 150$ words) is prepared. These descriptions:

\begin{itemize}
    \item Describe system setup, institutional structure, and task intent.
    \item Avoid failure diagnosis language.
    \item Avoid attributing blame, root cause, or governance breakdown.
    \item Avoid vocabulary specific to any analytical framework.
\end{itemize}

The purpose of these neutral descriptions is to serve as identical input seeds across Baseline, Protocol, and Independent generation conditions.

\paragraph{Freeze Procedure.}
After preparation, the following artifacts are frozen:

\begin{itemize}
    \item Case list (C01--C25)
    \item Case ID mapping
    \item All 25 neutral case descriptions
\end{itemize}

The frozen bundle is archived with:

\begin{itemize}
    \item Timestamp (UTC)
    \item Cryptographic hash (SHA-256)
    \item Version identifier (v1.0)
\end{itemize}

Once frozen, no wording edits, additions, or structural changes are permitted. Any modification requires incrementing the version number and re-hashing the full bundle.

This freeze ensures that:
\begin{enumerate}
    \item All generation conditions operate from identical inputs.
    \item No post hoc framing changes are introduced.
    \item Case sampling remains fixed and auditable.
\end{enumerate}